\paragraph{C.3av (Refuted, with specified counterexamples): the full low-dimensional ordinary-surjectivity range in positive characteristic.}
For every prime $p$ and every pair of integers
\[
                   r\ge1,\qquad r+2\le N\le2r+1,
\]
there is a specified algebraically closed field $k$ of characteristic
$p$, a regular excellent strictly henselian local $k$-ring $R$ of
dimension $N$, and a principal divisor with regular puncture for which
\[
 \mathrm{CH}^r(X)=0,
 \qquad \mathrm{CH}^r(Y)_{\mathrm{alg}}\otimes\mathbb Q\ne0.
\]
The coefficient field is a specified algebraic closure of a function
field over $\overline{\mathbb F}_p$, not the algebraic closure of a
finite field itself. Thus this result does not contradict the vanishing
of rational ordinary algebraically trivial cycles over the latter field.

\begin{proof}
\emph{Step 1: an elementary ordinary-top-cohomology hypersurface.}
Write
\[
 d=2r-N+2\ge1,\qquad s=N-r-2\ge0,
 \qquad r=d+s,\qquad N=d+2s+2,
 \qquad m=d+2.
\]
Choose the degree
\[
 e=\begin{cases}m,&p\nmid m,\\m+1,&p\mid m.\end{cases}
 \qquad\text{so that }p\nmid e.
\]
With variables $y_0,\ldots,y_{m-1}$ define
\[
 G=\begin{cases}
       \prod_i y_i,&e=m,\\
       (\prod_i y_i)(\sum_i y_i),&e=m+1,
    \end{cases}
 \qquad
 F_\mu=\sum_i y_i^e+\mu G.
\]
The hypersurface at $\mu=0$ is smooth by its Fermat gradient.
The same proper-singular-locus argument as in C.3at shows that
only finitely many parameters fail smoothness.

We give an elementary determinant calculation proving that some smooth
member has bijective Frobenius on $H^d(\mathcal O)$. The connecting
Cech presentation has basis $y^{-u}$ with $u_i\ge1$ and $\sum u_i=e$.
There is one basis vector $u=\mathbf1$ when $e=m$, and there are
$m$ vectors $u_j=\mathbf1+\epsilon_j$ when $e=m+1$.
Here $\epsilon_j$ is the $j$th coordinate vector. In these bases the
Hasse matrix is
\[
 A_{ij}(\mu)=[y^{p u_j-u_i}]F_\mu^{p-1}.
\]
This follows directly by multiplying the $p$th power of a Cech
representative by $F_\mu^{p-1}$, as in C.3at; it is also
\citet[Corollary~2.3, author p.~8]{D_HuangLianYauYuHasse}.
For $e=m$ its leading coefficient as a polynomial in $\mu$ is one.
For $e=m+1$, the coefficient of $\mu^{p-1}$ in $A_{ij}$ is the
coefficient of $y^{p\epsilon_j-\epsilon_i}$ in
$(\sum y_i)^{p-1}$. This is zero if $i\ne j$, because one exponent
is negative, and is one if $i=j$. Thus, putting
$g=1$ or $m$ respectively,
\[
                 \det A(\mu)=\mu^{g(p-1)}+
                                    \text{lower-degree terms}.
\]
There is therefore a nonempty open parameter set on which both
smoothness and $\det A\ne0$ hold, over every prime field.

Specify a parameter by the finite-field enumeration of C.3at: order
monic irreducible polynomials over $\mathbb F_p$ by degree and
lexicographic coefficient vector, put $\mu=z\bmod g_0(z)$ in the
resulting finite field, and take the first whose determinant is
nonzero and whose projective Jacobian ideals are all unit ideals.
This rule terminates by the nonempty-open argument. Denote the
resulting field by $\mathbb F_q$, and put
\[
 T_0=V(F_\mu)\subset\mathbb P^{d+1}_{\mathbb F_q}.
\]
It is smooth and geometrically integral, with
$H^i(\mathcal O)=0$ for $0<i<d$ and
$\dim H^d(\mathcal O)=g>0$. Frobenius on the latter group is
bijective by the determinant calculation.

\emph{Step 2: the Witt-vector conclusion, including the curve endpoint.}
Let $M_n=H^d(T_0,W_n\mathcal O)$ and $W=W(\mathbb F_q)$.
For $d\ge2$ the proof of C.3at gives
\[
 0\longrightarrow M_{n-1}\xrightarrow{V_n}M_n
                  \longrightarrow H^d(T_0,\mathcal O)\longrightarrow0.
\]
For $d=1$ the potential preceding connecting map is also zero:
\[
 H^0(T_0,W_n\mathcal O)=W_n(\mathbb F_q)
       \longrightarrow H^0(T_0,\mathcal O)=\mathbb F_q
\]
is surjective. This equality of global sections follows directly from
the Witt coordinates and $H^0(T_0,\mathcal O)=\mathbb F_q$.
Hence the same short exact sequence holds in the curve case.
In every dimension, the succeeding degree-$d+1$ group vanishes by
the finite Witt filtration and coherent cohomological dimension.

Write $F_n$ for coordinatewise Frobenius, $R_n$ for truncation,
and $\sigma$ for Frobenius on $W$. As in C.3at, $F_n$ is
$\sigma$-semilinear and bijective, $V_n$ is
$\sigma^{-1}$-semilinear, $R_n$ is surjective, and
\[
 p\operatorname{id}_{M_n}=V_nF_{n-1}R_n,
 \quad\operatorname{length}_W(M_n)=gn,
 \quad M_n/pM_n\simeq\mathbb F_q^g.
\]
A minimal set of $g$ generators gives a surjection $W_n^g\to M_n$;
the equal finite lengths make it an isomorphism. Lift compatible bases
through the surjective truncations. The proper inverse-limit comparison
in \citet[Section~2, author pp.~6--7]{D_BerthelotBlochEsnault2007}
then gives
\[
 H^d(T_0,W\mathcal O)\simeq W^g,
 \qquad
 E=H^d_{\mathrm{cris}}(T_0/W)[1/p]_{<1}\ne0,
\]
and every slope of $E$ is zero. The second assertion uses their
smooth proper comparison (1.1), equivalently Theorem~1.1. This proves
only the top coherent/unit-root statement required here; full ordinarity
in other degrees is unnecessary.

\emph{Step 3: a smooth suspension and an integral correspondence inverse.}
Over $k_0=\overline{\mathbb F}_p$ define
\[
 V_0=V\left(F_\mu(y)+\sum_{j=0}^{s-1}(a_j^e+b_j^e)\right)
                 \subset\mathbb P^{d+2s+1},
 \qquad n=\dim V_0=d+2s.
\]
For $s=0$ the paired coordinates are omitted. This hypersurface is
smooth. A simultaneous zero of its paired partials has every
$a_j=b_j=0$, since $p\nmid e$. If all remaining partials vanish,
Euler's identity $eF_\mu=\sum y_i\partial_iF_\mu$ gives
$F_\mu=0$, contradicting smoothness of $T_0$ unless all $y_i=0$.
Thus there is no projective singular point. The condition $p\nmid e$
is essential in this verification and was built into the construction.

Choose two distinct roots $\lambda_0,\lambda_1\in k_0$ of
$\lambda^e=-1$; they exist because $e\ge3$ and $p\nmid e$.
For $s>0$ the equations $a_j=\lambda_i b_j$ define the join
of $T_0$ with a linear $\mathbb P^{s-1}$. Its incidence resolution
$Z_i\subset T_0\times V_0$ is a $\mathbb P^s$-bundle over $T_0$,
with fibre coordinates
\[
 [\lambda_i b_0:b_0:\cdots:\lambda_i b_{s-1}:b_{s-1}:
                                  t y_0:\cdots:t y_{d+1}].
\]
It defines $\Gamma_i=[Z_i]\in\mathrm{CH}^{r}(T_0\times V_0)$.
When $s=0$ put $\Gamma_i=\Delta_{T_0}$.
The two different linear equations in each pair force $a_j=b_j=0$;
the two joins meet exactly along $T_0$ and avoid their vertices there.
Their tangent directions in the $2s$ paired coordinates are
complementary, because $\lambda_0\ne\lambda_1$. Their intersection
is transverse with multiplicity one. The correspondence composition
therefore gives, integrally,
\[
                    {}^t\Gamma_1\circ\Gamma_0=\Delta_{T_0}.
\]
This is the geometric calculation of C.3ai and uses no
characteristic-zero assumption.

\emph{Step 4: nonsupport by the hyperplane class, with exact slope shifts.}
Set $K=k_0(T_0)$, $k=\overline K$, and write $T,V$ for base change
to $k$. Let $P\in T(k)$ be the geometric generic point and choose
a constant point $o\in T_0(k_0)$. Both are fully specified by the
function-field inclusion and, for $o$, finite-field point enumeration.
The difference $\alpha=[P]-[o]$ is ordinarily algebraically trivial:
use $T$ itself when $d=1$, and a smooth connected sufficiently ample
complete-intersection curve through the two points when $d>1$.
The latter existence and actual fixed-product family are as in C.3at.
Consequently $z=(\Gamma_0)_*\alpha$ is an actual ordinarily
algebraically trivial class of codimension $r$ on $V$.

Let $H=c_1(\mathcal O_V(1))$. Suppose $z=H\beta$ in rational
Chow groups. The relation descends to a finite extension of $K$.
Taking its norm, spreading the generic first factor, and closing all
cycles gives
\[
 b\Gamma_0=H B+
       b\bigl(T_0\times(\Gamma_0)_*[o]\bigr)+\Theta,
 \qquad b\ne0,
\]
in $\mathrm{CH}^r(T_0\times V_0)_{\mathbb Q}$, where
$B\in\mathrm{CH}^{r-1}(T_0\times V_0)_{\mathbb Q}$ and
$\Theta$ is supported on $Z\times V_0$ for a proper divisor
$Z\subset T_0$. The constant correspondence is retained explicitly.
All these finite data descend to a finite extension of $\mathbb F_q$.
We make one such extension also containing the two roots
$\lambda_i$ and use Frobenius-equivariant rational crystalline
correspondences over that field.

The nonzero slope-zero space $E\subset H^d_{\mathrm{cris}}(T_0)[1/p]$
maps by $\Gamma_0$ into
\[
                       H^n_{\mathrm{cris}}(V_0)[1/p](s).
\]
The transpose correspondence is its left inverse, by the integral
composition identity. Thus this map is injective. The constant term
acts as zero in degree $d>0$. The term $\Theta$ also kills $E$:
alter each source-support image to a smooth projective variety of
dimension at most $d-1$, lift the correspondence rationally, and use
\[
 H^d_{\mathrm{cris}}(Z')[1/p]_{<1}
         =H^d(Z',W\mathcal O)[1/p]=0.
\]
The rational lifting, finite-field descent and source-support argument
are proved in Step~3 of C.3at, using de Jong's Theorem~4.1 and
the Berthelot--Bloch--Esnault comparison; they apply unchanged here.

It remains to show that $H B$ kills $E$. With the convention that
$(1)$ divides Frobenius eigenvalues by $q$, the action of $B$ is
\[
       E\longrightarrow
       H^{d+2s-2}_{\mathrm{cris}}(V_0)[1/p](s-1)
         =H^{n-2}_{\mathrm{cris}}(V_0)[1/p](s-1).
\]
If $s=0$, this target has all slopes at least one: smooth proper
crystalline cohomology has nonnegative slopes and the twist is $(-1)$.
Negative cohomological degrees are zero. Hence the map is zero.

If $s\ge1$, rational crystalline weak Lefschetz gives
\[
 H^{n-2}_{\mathrm{cris}}(\mathbb P^{n+1})[1/p]
        \simeq H^{n-2}_{\mathrm{cris}}(V_0)[1/p].
\]
Here is the exact needed proof, avoiding any integral weak-Lefschetz
claim. The complement $U=\mathbb P^{n+1}\setminus V_0$ is smooth
affine of dimension $n+1$. Rigid cohomology agrees with
Monsky--Washnitzer cohomology and vanishes above that dimension.
Rigid Poincare duality therefore makes $H^j_{\mathrm{rig},c}(U)=0$
for $j<n+1$. These are the inputs explicitly recalled in
\citet[proof of Corollary~1.2, author p.~4]{D_BerthelotBlochEsnault2007}.
The compact-support localization sequence yields the displayed
isomorphism, indeed an isomorphism in every degree below $n$.
All maps commute with Frobenius. If $n-2$ is odd, this group is zero;
otherwise its sole untwisted slope is $(n-2)/2$. After the twist
$(s-1)$ the sole slope is
\[
                         \frac{n-2}{2}-(s-1)=\frac d2>0.
\]
There is no Frobenius-equivariant map from a slope-zero isocrystal to
this target. Thus $B$, and consequently $HB$, kills $E$ in every case.
The spread equality would contradict injectivity of $b\Gamma_0$ on
$E$. We conclude
\[
                 z\ne0\quad\text{in}\quad
       \mathrm{CH}^r(V)_{\mathbb Q}/H\mathrm{CH}^{r-1}(V)_{\mathbb Q}.
\]

\emph{Step 5: the actual strict local example.}
Specify
\[
 \begin{aligned}
 R&=k[a_0,b_0,\ldots,a_{s-1},b_{s-1},y_0,\ldots,y_{d+1}]
                        _{(a,b,y)}^{\mathrm{sh}},\\
 f&=F_\mu(y)+\sum_{j=0}^{s-1}(a_j^e+b_j^e).
 \end{aligned}
\]
The ambient dimension is $d+2s+2=N$. The hypersurface cone is
normal and regular away from its vertex, by the gradient check and
the hypersurface depth criterion. Its homogeneous ring is the full
section ring of $\mathcal O_V(1)$, generated in degree one.
The quotient $R/(f)$ is the strict henselization of the cone's local
ring, so $Y$ is its regular puncture. C.3e injects the nonzero quotient
class $z$ into $\mathrm{CH}^r(Y)_{\mathbb Q}$. The incidence
correspondence and the curve family defining $\alpha$ pull back to
actual families on the fixed puncture, so the integral class lies in
$\mathrm{CH}^r(Y)_{\mathrm{alg}}$.

The ambient $R$ is regular noetherian and contains the infinite field
$k$. Kerz's integral Milnor Bloch formula, Theorem~7.5, together with
exactness of global sections on the local Zariski space, gives
$\mathrm{CH}^r(\operatorname{Spec}R)=0$ just as in C.3at.
Chow localization gives $\mathrm{CH}^r(X)=0$. This proves the
claimed rational failure throughout the stated interval.
\end{proof}

The proof does not assert failure at $N\ge2r+2$. Its top-cohomology
argument uses rational crystalline slopes, not positive-characteristic
regular differential forms, and all transfer degrees are inverted in
the characteristic-zero cohomological coefficient field. The actual
cycle lives over the displayed transcendental algebraically closed
extension; over $\overline{\mathbb F}_p$ the rational ordinary group
still vanishes.
